\documentclass[11pt]{article}
\usepackage[margin=1in]{geometry}
\usepackage[T1]{fontenc}
\usepackage[utf8]{inputenc}
\usepackage{lmodern}
\usepackage{amsmath,amssymb,mathtools,booktabs,graphicx,natbib,hyperref}
\usepackage[nameinlink,noabbrev]{cleveref}

\title{Numerical Selection of Forecast-Optimal Smoothness and Financial Trend Recurrence}
\author{Heriberto Espino Montelongo}
\date{\today}
\newcommand{\R}{\mathbb{R}}

\begin{document}
\maketitle

\begin{abstract}
We study numerical selection of normalized smoothness for finite-difference
penalized trend estimation when the target is chronological forecast loss.
Rather than searching an arbitrary bounded range of the penalty parameter, we
parameterize the problem by a normalized smoothness index on the compact domain
$[0,1]$. Because the forecast-validation objective may have multiple local
minima, the numerical problem is to isolate and refine its stationary points
without evaluating a dense uniform grid. The selected trend is then used as a
forecast reference path for studying first-return and crossing times in
financial price series. This manuscript is a research skeleton; empirical
claims are deferred until the numerical method and recurrence protocol are
frozen and validated.
\end{abstract}

\section{Penalized trend estimation}

For $y\in\R^N$, let $D_d$ be the order-$d$ finite-difference matrix and
$Q=D_d^\top D_d$. Then
\[
\widehat\tau_\lambda
=
\arg\min_{\tau\in\R^N}
\left\{
\|y-\tau\|_2^2+\lambda\|D_d\tau\|_2^2
\right\}
=
(I+\lambda Q)^{-1}y.
\]

\section{Normalized smoothness}

Let $\delta_1,\ldots,\delta_{N-d}$ be the positive eigenvalues of $Q$.
Define
\[
S(\lambda)
=
1-
\frac{1}{N-d}
\sum_{j=1}^{N-d}
\frac{1}{1+\lambda\delta_j}.
\]
Then $S(0)=0$, $S(\lambda)\to1$ as $\lambda\to\infty$, and
\[
S'(\lambda)
=
\frac{1}{N-d}
\sum_{j=1}^{N-d}
\frac{\delta_j}{(1+\lambda\delta_j)^2}
>0.
\]

\section{Forecast-optimal smoothness}

For fixed forecasting protocol and horizon $h$, let $CV_h(S)$ be pooled
chronological forecast loss. The principal object is
\[
\boxed{
S_h^\star=\arg\min_{S\in[0,1]}CV_h(S).
}
\]

If $f(\lambda)$ is the same objective in penalty coordinates and
$F(S)=f(\lambda(S))$, then
\[
F'(S)=\frac{f'(\lambda)}{S'(\lambda)}.
\]
Thus interior stationary points are preserved by the reparameterization.

\section{Multiple-local-minimum numerical search}

The proposed method starts from a small partition of the smoothness domain,
adaptively refines intervals with evidence of stationary structure, brackets
derivative sign changes, applies Brent refinement, classifies stationary
points, and compares all detected local minima with exact boundaries. A dense
GPU smoothness grid is retained as a numerical reference.

\section{Financial trend recurrence}

At forecast origin $T$, only information through $T$ may determine the selected
smoothness and trend. Let $\widehat\tau_{T+k\mid T}$ be the frozen forecast
trend path and let $x_t=\log P_t$. Define
\[
g_{T,k}=x_{T+k}-\widehat\tau_{T+k\mid T}.
\]
A crossing recurrence time is
\[
R_T^{\mathrm{cross}}
=
\inf\left\{
k\ge1:g_{T,k}g_{T,0}\le0
\right\}.
\]
A tolerance-band recurrence is defined analogously. Events that do not recur
inside a pre-specified horizon are right-censored.

\section{Experimental design}

First validate the numerical optimizer against dense reference grids. Then
freeze the forecasting and recurrence protocol before evaluating broad
ETFs/indices, held-out equities, and cryptocurrency. The main paper does not
jointly adapt difference order and window length.

\section{Limitations}

Recurrence to a forecast trend is not equivalent to stationarity or classical
mean reversion. Forecast accuracy and recurrence statistics do not imply
trading profitability. Dependence induced by overlapping forecast origins must
be addressed before inferential claims.

\bibliographystyle{plainnat}
\IfFileExists{literature/references.bib}{
  \bibliography{literature/references}
}{
  \bibliography{../literature/references}
}
\end{document}
